\begin{table}[tb] % Inizio ambiente tabella flottante; [tb] = top o bottom della pagina
\centering % Centra la tabella
\caption{Esempio di tabella con vari campi, spunte e croci.} % Didascalia
\begin{adjustbox}{width=0.9\textwidth} % Adatta la tabella al 90% della larghezza del testo
{\normalsize % Imposta la dimensione normale del font
\setlength{\tabcolsep}{8pt} % Imposta la spaziatura orizzontale tra le colonne

\begin{tabular}{cccrrrr} % Crea una tabella con 7 colonne:
                         % c = allineamento centrato, r = allineamento a destra
\toprule % Riga orizzontale superiore spessa (booktabs)

% \makecell permette di inserire testo formattato o su più righe all’interno di una cella.
% È utile soprattutto nelle intestazioni per evitare l’uso di \parbox o \multicolumn.
% Qui viene usato per centrare e rendere in grassetto il testo delle intestazioni.
\makecell[c]{\textbf{Frutto}} & 
\makecell[c]{\textbf{Arancione}} & 
\makecell[c]{\textbf{Si sbuccia}} & 
\makecell[c]{\textbf{Peso}} &   
\makecell[c]{\textbf{Prezzo}} &   
\makecell[c]{\textbf{Zucch.}} &   
\makecell[c]{\textbf{Voto (1–10)}} \\  
% Il simbolo & separa le colonne all'interno di una riga.
% Ogni cella di una riga è delimitata da un &.
% Il comando \\ indica la fine della riga e l’inizio di una nuova.
\midrule % Riga orizzontale media

Arancia & \checkmark  & \checkmark & 0.200 & 2.50 & 9.0 & 8 \\ % Esempio di riga: & separa le colonne, \\ chiude la riga
Mandarino & \checkmark & \checkmark & 0.120 & 3.20 & 8.5 & 9 \\
Limone & \ding{55} & \checkmark & 0.100 & 3.50 & 2.5 & 6 \\
\midrule

Mela & \ding{55} & \ding{55} & 0.180 & 2.80 & 10.0 & 8 \\
Banana & \ding{55} & \checkmark & 0.150 & 1.10 & 12.2 & 9 \\
Kiwi & \ding{55} & \checkmark & 0.120 & 4.00 & 9.1 & 7 \\
Mango & \checkmark & \checkmark & 0.300 & 5.00 & 14.0 & 10 \\
\bottomrule % Riga inferiore spessa

\end{tabular} % Fine ambiente tabular
} % Fine del blocco \normalsize
\end{adjustbox} % Fine dell’adattamento di larghezza
\label{tab:table_example} % Etichetta per riferimenti
\vspace{-10pt} % Riduce lo spazio sotto la tabella
\end{table} % Fine ambiente tabella
